\section{Explicit height functoriality and complete arithmetic dynamics}
\label{sec:dynamics}

A lower height bound must account for common factors removed from the image
coordinates. Bounding the displayed coordinates alone would not suffice.
The following determinant construction controls those factors and supplies
an effective finiteness theorem. Its classical setting is the height theory
of arithmetic dynamics \cite{silverman,silverman2011}; the proof here is
entirely in integer arithmetic.

\subsection{Two finite polynomial certificates}

Let
\[
 F(X,Y)=\sum_{i=0}^d f_iX^{d-i}Y^i,\qquad
 G(X,Y)=\sum_{i=0}^d g_iX^{d-i}Y^i,\qquad d\ge1.
\]
For a nonnegative integer $r$, a homogeneous binary form of degree $r$ is
represented by its $r+1$ coefficient list. Multiplication is finite
convolution. Form the $2d\times2d$ integer matrix $S$ of the map
\[
   (A,B)\longmapsto AF+BG,
\]
where $A,B$ have degree $d-1$ and the output has degree $2d-1$.
Use the ordered monomial basis $X^r,X^{r-1}Y,\ldots,Y^r$ in every space.
Let $R=\det S$. The nondegeneracy hypothesis for this section is the
decidable arithmetic condition $R\ne0$. This is the homogeneous resultant
condition, defined here by its matrix rather than by an algebraic-closure
criterion.

\begin{lemma}[Computed Bezout certificates]\label{lem:resultant}
The first and last columns of $\operatorname{adj}S$ provide forms
$A_X,B_X,A_Y,B_Y$ of degree $d-1$ with
\begin{align}
 A_XF+B_XG&=R X^{2d-1},\label{eq:bezout-X}\\
 A_YF+B_YG&=R Y^{2d-1}.\label{eq:bezout-Y}
\end{align}
These identities are provable in $\iPRA$ and are checked coefficientwise.
If
\[
 M=\max(1,|f_0|,\ldots,|f_d|,|g_0|,\ldots,|g_d|),
 \quad K=(2d-1)!M^{2d-1},
\]
every coefficient of the four certificate forms has absolute value at most
$K$.
\end{lemma}
\begin{proof}
The adjugate identity gives $S\operatorname{adj}S=RI$; its first and last
columns are precisely the two coefficient identities. Each cofactor is a
sum of $(2d-1)!$ signed products of $2d-1$ entries of $S$. Every entry of
$S$ is a coefficient of $F$ or $G$, or zero. Bounding each product by
$M^{2d-1}$ proves the claim. Determinants, permutations, signs and coefficient
convolution are finite sums and products with the displayed bounds; the
finite proof of the adjugate identity was given in
Section~\ref{sec:heights}.
\end{proof}

Set
\[
 L=\max\bigl(\|A_X\|_1+\|B_X\|_1,
             \|A_Y\|_1+\|B_Y\|_1\bigr),\qquad U=(d+1)M.
\]
Since $R\ne0$, $L\ge1$. The coarse bound
\begin{equation}\label{eq:coarse-L}
                L\le 2d(2d-1)!M^{2d-1}
\end{equation}
is useful if one does not wish to keep the smaller computed constant.

\begin{theorem}[Two-sided height functoriality on $\PP^1$]
\label{thm:height-map}
If $R\ne0$, the rule $f([x:y])=[F(x,y):G(x,y)]$ is a total map on
$\PP^1(\Q)$. For every rational projective point $P$,
\begin{equation}\label{eq:height-map}
              \Ht(P)^d\le L\Ht(f(P)),\qquad
              \Ht(f(P))\le U\Ht(P)^d.
\end{equation}
The map, constants, and proofs are uniform primitive recursive constructions
from the finite coefficient lists.
\end{theorem}
\begin{proof}
Take a primitive integral pair $(x,y)$ and put $H=\max(|x|,|y|)\ge1$.
If $F(x,y)=G(x,y)=0$, the two identities imply
$R x^{2d-1}=R y^{2d-1}=0$. Since $R\ne0$ and $2d-1>0$, this would give
$x=y=0$. Thus the image is nonzero. Homogeneity establishes that its
projective class is independent of the chosen representative.

Put $T=\max(|F(x,y)|,|G(x,y)|)$ and
$g=\gcd(|F(x,y)|,|G(x,y)|)\ge1$. The identities show that $g$ divides
both $R x^{2d-1}$ and $R y^{2d-1}$. The two powers are coprime; multiplying
their Bezout identity by $R$ shows $g\mid R$. Consequently $g\le|R|$.
This is the required control of projective cancellation.

Choose a coordinate of absolute value $H$ and use the corresponding
certificate. Each coefficient form has degree $d-1$, so
\[
 |R|H^{2d-1}\le L H^{d-1}T.
\]
Cancel the positive integer $H^{d-1}$ to obtain
$|R|H^d\le LT$. Since $T=g\Ht(f(P))$ and $g\le|R|$, we have
\[
 |R|H^d\le Lg\Ht(f(P))\le L|R|\Ht(f(P)).
\]
Cancellation of $|R|>0$ proves the lower bound. For the upper bound,
$T\le(d+1)MH^d=UH^d$, and division by $g$ cannot increase height.
All case distinctions here concern integer comparisons or finite lists.
\end{proof}

The proof also works with any supplied forms satisfying
\eqref{eq:bezout-X}--\eqref{eq:bezout-Y} and any nonzero $R$; it need not
recompute the adjugate. The same two finite identities prove nonvanishing,
bound the gcd, and give the lower height estimate. This is a useful
certificate interface for formalization.

\subsection{A global finiteness theorem with an actual decision procedure}

Define the iterate by $f^0(P)=P$ and $f^{n+1}(P)=f(f^n(P))$. A point is
\emph{preperiodic} if
\[
             \exists i\,\exists j\ (i<j\ \wedge\ f^i(P)=f^j(P)).
\]
This is initially an unbounded existential statement. The next theorem
reduces it to a bounded computation.

\begin{theorem}[Primitive recursive Northcott theorem for rational dynamics]
\label{thm:preperiodic}
Suppose $d\ge2$ and $R\ne0$. Every rational preperiodic point of $f$ has
height at most $L$. There is a primitive recursive algorithm that returns
\emph{all} rational preperiodic points, together with a repetition witness
for each, and whose correctness is provable in $\iPRA$.
In particular there are at most $(2L+1)^2$ such points.
\end{theorem}
\begin{proof}
If $H=\Ht(P)>L$, then
$H^d\ge H^2>LH$. By \eqref{eq:height-map}, $L\Ht(f(P))\ge H^d>LH$,
so $\Ht(f(P))>H$. Induction on a numerical iteration counter therefore
shows that all subsequent heights strictly increase. Given $i<j$, a
finite induction from $i$ through $j$ proves
$\Ht(f^i(P))<\Ht(f^j(P))$. Such a point cannot have a repetition witness.
Thus any preperiodic point has height at most $L$.

Every iterate of a preperiodic point is preperiodic: apply $f^t$ to a
witness equality and use $f^i(f^t(P))=f^t(f^i(P))$, an identity proved by
finite recursion. It follows that its entire forward orbit remains in
the explicit finite list
\[
             S=\mathcal L(1,L),\qquad N=|S|\le(2L+1)^2.
\]

The algorithm first constructs $S$. For each $P\in S$, start an empty
visited list and follow its orbit for at most $N+1$ inspected states. If
the current height exceeds $L$, reject $P$. If the current state equals
a previously visited state, retain $P$ and the two visit indices. Otherwise
append it and continue. The transition can be made a map on the bounded
set $S$ plus one escape flag: after escape, keep the flag unchanged.

No candidate can remain unclassified. If $N+1$ states were inspected
without escape or repetition, they would be $N+1$ different members of an
$N$-element list. Here is a finite arithmetic proof of the contradiction.
Find each state's index in $S$ by bounded search. For every $r<N$, let
$c_r$ count the occurrences of that index among the $N+1$ visits. Absence
of repetitions gives $c_r\le1$, hence $\sum_{r<N}c_r\le N$. On the
other hand each visit has exactly one index, so summing the same incidence
table by visits gives $\sum_{r<N}c_r=N+1$. Both equalities are bounded
double-sum identities.

Soundness of retention is the returned equality with two distinct indices.
If a candidate is preperiodic, all its iterates lie in $S$, so rejection
by escape cannot occur; the preceding pigeonhole argument then retains it.
The earlier height bound proves that no preperiodic point was omitted from
the candidate list.

The certificate construction, construction of $S$, and the loop counts are
bounded by primitive recursive functions of the coefficient lists. Before
escape, raw image coordinates are at most $UL^d$. Keeping the escape flag
prevents uncontrolled growth of the stored states. Thus the whole procedure
is a fixed primitive recursive algorithm, not a merely terminating search
whose running time has been left unspecified. Its invariants are Boolean
tests, list membership, and bounded counting equations.
\end{proof}

One may use the coarser right side of \eqref{eq:coarse-L} everywhere in
place of $L$ and still obtain a primitive recursive bound depending only
on $d,M$. The bound is usually far too large for practical computation.
Primitive recursive totality and computational efficiency are separate issues.

\subsection{Two complete examples}

For $f(t)=t^2$, use $F=X^2$, $G=Y^2$ and the identities
$XF=X^3$, $YG=Y^3$. Here $R=1$ and $L=1$.
The rational projective points of height at most one are
$\infty,0,1,-1$. Their images are respectively $\infty,0,1,1$.
All four are preperiodic, and Theorem~\ref{thm:preperiodic} proves that
there are no others.

For $f(t)=t^2-1$, use $F=X^2-Y^2$, $G=Y^2$ and
\[
                 XF+XG=X^3,\qquad YG=Y^3.
\]
Thus $R=1$ and $L=2$. The complete height-two list is
\[
             \infty,\ 0,\ \pm1,\ \pm2,\ \pm\tfrac12.
\]
The four points $\infty,0,1,-1$ are retained: $\infty$ is fixed,
$0\leftrightarrow-1$ is a cycle, and $1\mapsto0$.
Both $2$ and $-2$ map to $3$, whose height exceeds two. Both
$1/2$ and $-1/2$ map to $-3/4$, whose height exceeds two.
They cannot be preperiodic by the proven escape bound. Consequently
\[
          \operatorname{PrePer}(t^2-1,\Q)=\{\infty,0,1,-1\}.
\]
These calculations are reproduced by the exact Python companion. Their
completeness comes from the height proof, not from a numerical search cutoff.
